\section{Implementation and Evaluation}
\label{sec:eval}

\paragraph{Setup.} The extractor is trained on SOCOFing~\cite{Shehu18} (6,000 fingerprints, 600
subjects $\times$ 10 fingers) at $224\times224$, with an identity-disjoint split: 4,800 training
images and 1,200 held-out identities that appear in no training or calibration step. Enrolment
uses the \texttt{Real} captures and probes use the \texttt{Altered-*} sets. Confidence intervals
throughout are 95\% bootstrap over \emph{identities}, not pairs --- pairs sharing a finger are
not independent, and resampling pairs understates the interval. Hardware: RTX 4050 (6\,GB); the
FLPSI layer runs against a reference implementation of~\cite{BuiCong25} instrumented to count
protocol bytes.

\subsection{Accuracy}

\begin{table}[htbp]
\centering
\caption{EER with 95\% identity-level bootstrap intervals, at a matched bit budget. Collisions
count impostor pairs at $H = 0$ out of ${\approx}1.44$M.}
\label{tab:acc}
\small
\begin{tabular}{@{}lcccr@{}}
\toprule
Coder & Altered-Easy & Altered-Medium & Altered-Hard & $H{=}0$ \\
\midrule
Float embedding (ceiling) & 0.33\% & --- & --- & --- \\
\textbf{Super-Bit 128} & \textbf{1.88\%} [1.56, 2.20] & 4.48\% [3.90, 4.94] &
  7.04\% [6.31, 7.45] & \textbf{2} \\
ITQ 128 & 4.55\% [3.91, 5.31] & 7.11\% [6.39, 7.73] & 9.06\% [8.47, 9.75] & 517 \\
ITQ 16 & 6.13\% [5.54, 6.80] & 8.49\% [7.84, 9.19] & 11.58\% [10.56, 12.20] & 6{,}394 \\
Naive sign 16 & 10.12\% & --- & --- & --- \\
\bottomrule
\end{tabular}
\end{table}

Binarisation costs $+1.55$ percentage points of EER against the $0.33\%$ floating-point ceiling
(Table~\ref{tab:acc}). Two comparisons matter more than the headline.

First, \textbf{the bit budget}. Comparing a 128-bit Super-Bit code against a 16-bit ITQ code
would be a mismatch, so ITQ is fitted at 128 bits as well. Super-Bit still wins by more than
$2\times$ ($1.88\%$ vs $4.55\%$), so its advantage comes from block orthonormalisation rather
than from a larger budget. It also wins by $250\times$ on code collisions, which
\S\ref{sec:floor} shows is the metric that ultimately binds.

Second, \textbf{degradation}. All three coders degrade gracefully and monotonically across the
three alteration levels, with Super-Bit's ordering preserved throughout.

We note against Blind-Touch's reported $0.7\%$ EER that our interval $[1.56, 2.20]$ does not
contain it; the comparison is also across different training regimes and should be read as
indicative rather than a controlled contrast.

\paragraph{Whitening ablation.} Rescaling embedding dimensions by the head's final-layer weights
before projection helps an under-trained model but \emph{hurts} the well-trained one
($1.88\% \to 2.98\%$), because a well-conditioned embedding gains nothing from reweighting and
loses from the added variance. The recommended configuration omits it.

\paragraph{Code-length ablation.} At 64/128/256 bits the EER is $2.38\%$ / $1.88\%$ / $1.89\%$
with storage 8/16/32\,B per user. Accuracy saturates at 128 bits: this is the 16-dimensional
information ceiling, not a coding limitation, and it is the same ceiling that produces the
distributional problem of \S\ref{sec:params}.

\subsection{Operating point}

The accuracy figures above describe the \emph{representation}. The deployed decision rule is the
FLPSI accept model, and \S\ref{sec:params} shows that at the published parameters its per-record
false-accept rate is $4.578\times10^{-2}$ --- $1{,}550\times$ the value predicted under the
uniform-code model --- with a query-level rate that exceeds $50\%$ by $N = 100$. We do not repeat
that analysis here; we note only that \textbf{an EER measured on the code is not the error rate
of the deployed system}, and that the two must be reported together.

\subsection{Communication}

Online per-query communication is linear in the database,
$\text{comm} \approx 0.25\,\mathrm{MB} + 1.54\,\mathrm{KB}\cdot N$, measured on the query channel
in both directions with garbled-circuit transfer excluded as one-time, data-independent
preprocessing.

\begin{center}\small
\begin{tabular}{@{}lrrrrrr@{}}
\toprule
$N$ & 500 & 1{,}000 & 5{,}000 & 10{,}000 & 50{,}000 & $10^{6}$ (extrap.) \\
\midrule
online comm & 1.00\,MB & 1.76\,MB & 7.77\,MB & 15.28\,MB & 75.40\,MB & ${\sim}1.5$\,GB \\
\bottomrule
\end{tabular}
\end{center}

The extrapolation matches Bui and Cong's reported 153\,MB at $10^{5}$ and 1{,}533\,MB at $10^{6}$,
confirming that the composition inherits the FLPSI communication profile unchanged --- the
quantiser adds nothing, it only fixes the 128-bit width.

\textbf{This is where the design loses.} Blind-Touch sends ${\sim}1.1$\,MB per query at constant
cost below 8,192 users; we exceed that beyond ${\sim}500$ records. Communication is not a
selling point of this work and we do not present it as one. The protocol runs in a constant
number of rounds over a field with $|\mathbb{F}| \ge 2^{\kappa + \log N}$, which for
$N = 5\times10^{4}$ and $\kappa = 40$ means a 56-bit field.

\subsection{Latency and storage}

\begin{center}\small
\begin{tabular}{@{}lrrl@{}}
\toprule
phase & $N = 1{,}000$ & $N = 5{,}000$ & scaling \\
\midrule
CNN forward (client, GPU) & 0.80\,ms & 0.80\,ms & $O(1)$ \\
Quantise (client) & 0.002\,ms & 0.002\,ms & $O(1)$ \\
FLPSI online & 187\,ms & 823\,ms & ${\sim}0.16$\,ms/record \\
\midrule
total query-time & ${\sim}188$\,ms & ${\sim}824$\,ms & \\
FLPSI offline (preprocessing) & 553\,ms & 1{,}946\,ms & ${\sim}0.38$\,ms/record \\
\bottomrule
\end{tabular}
\end{center}

The quantiser costs one $128\times16$ matrix multiply and a threshold compare --- $2.2\,\mu$s,
negligible against everything else. Server-side storage is \textbf{16\,B per user} against
Blind-Touch's 1.67\,KB, a $107\times$ reduction, with \textbf{no 117\,MB Galois key} and no
slot-packing ceiling. Storage and key material are where this design wins, and the win is
structural rather than a constant factor.
